\section{The Kingdom of Heaven}

\paragraph{Church of John}
We've mentioned the Church of John\footnote{\url{https://www.meditationsonthetarot.com/the-church-of-john}} before. There is a direct connection from John to
\textbf{Clement of Alexandria}:

\textbf{St John the Theologian} $\Rightarrow $ Irenaeus $\Rightarrow $ Justin Martyr $\Rightarrow $ Clement of
Alexandria

Subsequently, there is this intellectual chain:

\textbf{Clement of Alexandria} $\Rightarrow $ Origen $\Rightarrow $ Gregory of Nyssa\footnote{\url{https://gornahoor.net/?p=6071}} $\Rightarrow $ Maximos the
Confessor

More recently, \textbf{Hans Urs von Balthazar} and \textbf{Valentin Tomberg} are representatives of the Church of John.

In this tradition, there is the distinction made between exoteric and esoteric paths. The latter is reserved to those
capable of and interested in pursuing that path. The goal is gnosis. In researching this post, I listened to a youtube
lecture by a theologian who really had no understanding of gnosis. He thought is had something to do with passing on a
secret. If you recall, Tomberg explicitly denies that gnosis is a secret. That is because gnosis is personal knowledge
and cannot be shared. However, the means to achieve it can be shared.

This is a forgotten tradition and even a forbidden tradition. Recently a super-correct person informed us that we should
ignore the Fathers because they held all sorts of heretical ideas. Don't get sucked into that
vortex. Several years ago, I bookmarked several pages in Volume Two of the \emph{Philokalia}. I
don't recall why, but I am providing some of the texts from those pages. Note that images
involving space and time are not understood literally. Instead, these need to be understood as representing different
spiritual states. The next several sections are from \emph{Two Hundred Texts on Theology} by \textbf{St Maximos the
Confessor}, written in the second century.

\paragraph{Eternal Mansions}
Here we see that the mansions are not spatial locations but rather refer to spiritual qualities.

\begin{quotex}
How does the difference between the eternal dwelling places arise? Scholars propose two options:

\end{quotex}
\begin{enumerate}
\item \begin{quotex}
The difference is in their actual locality 
\end{quotex}
\item \begin{quotex}
The difference arises from our conception of spiritual quality and quantity peculiar to each mansion. 
\end{quotex}
\end{enumerate}
\begin{quotex}
Whoever knows the meaning of ``the kingdom of God is within you" and ``In my Father's house are many
mansions" will prefer the second explanation. [\#89] 

\end{quotex}
\paragraph{The Kingdom of Heaven}
This passage illustrates several ideas very densely:

\begin{itemize}
\item The inner essences of created things exist in God pre-eternally. 
\item By knowing those essences, we reach the kingdom of heaven. 
\end{itemize}
\begin{quotex}
The kingdom of heaven consists in possessing an inviolate and pre-eternal knowledge of created things through perceiving
their inner essences as they exist in God. [\#90] 

\end{quotex}
\paragraph{Time and the Kingdom of Heaven}
Once again, a dense passage.

\begin{itemize}
\item The arrival of the kingdom of heaven is not an event in time. 
\item The kingdom of heaven is not material since it cannot be observed as an object of sense. 
\item The distance from the kingdom of heaven reflects the inner state of the individual 
\end{itemize}
\begin{quotex}
The text, ``The kingdom of heaven has drawn near", does not imply any temporal limitation. For the kingdom does not come
in a way that can be observed: one cannot say, ``Look, it is here" or ``look, it is there". The phrase has reference to
the relationship which the saints have with the kingdom, each according to his inner state. For ``the kingdom of God is
within you." [\#91] 

\end{quotex}
\paragraph{The Kingdom of God}
This section needs to be unpacked:

\begin{itemize}
\item The kingdom of God needs to be actualized in each believer. 
\item That is achieved by clearing out all worldly distractions from one's inner life. 
\item Then, my false self dies and my real self lives in me 
\end{itemize}
\begin{quotex}
The kingdom of God the Father is present in all believers in potentiality; it is present in actuality in those who,
after totally expelling all natural life of soul and body from their inner state, have attained the life of the Spirit
alone and are able to say, ``I no longer live, but Christ lives in me." [\#92] 

\end{quotex}
\paragraph{Quality of Righteousness}
We see three options for understanding the kingdom of heaven. They are all correct, since they depend on the qualities
(i.e., what he has attained) of the person. From lower to higher, these are:

\begin{enumerate}
\item It is a prolongation of the human soul as it had existed on earth. 
\item The being reaches an angelic state as we described recently. 
\item The third state is for the reader to try to understand. 
\end{enumerate}
\begin{quotex}
Some say:

\end{quotex}
\begin{itemize}
\item \begin{quotex}
The kingdom of heaven is the way of life which the saints lead in heaven 
\end{quotex}
\item \begin{quotex}
It is a state similar to that of the angels, attained by those who are saved 
\end{quotex}
\item \begin{quotex}
It is the very form of the divine beauty of those who ``wear the image of Him who is from heaven". 
\end{quotex}
\end{itemize}
\begin{quotex}
In my judgment each of these three views is correct. For the grace of the kingdom is given to all according to the
quality and quantity of the righteousness that is in them. [\#93] 

\end{quotex}
\paragraph{Gnostic Philosophy}
Note here the method of biblical interpretation, since this is an example of the highest level. This is a commentary on
the incarnation and ascension of Jesus understood as spiritual states.

\begin{itemize}
\item The Logos incarnates in us as the commandments from the Father to the extent that we engage in holy warfare. 
\item After success in spiritual warfare, we move on to contemplation and reach gnosis. 
\item ``Return to the Father" means achieving a transcendent intellectual state. 
\end{itemize}
\begin{quotex}
So long as we are manfully engaged in the holy warfare of ascetic or practical philosophy we retain with us the Logos,
who in the form of the commandments came from the Father into this world. But when we are released from our ascetic
struggle with the passions and are declared victor over both them and the demons, we pass, by means of contemplation,
to gnostic philosophy; in this way we allow the Logos mystically to leave the world again and make his way to the
Father.

Hence it is that the Lord says to His disciples, ``You have loved me and have believed that I come from God. I came from
the Father and have come into the world; again I leave the world and make my way to the Father."

\begin{itemize}
\item\textbf{World}: this means the hard task of practicing the virtues. 
\item\textbf{Father}: this means the intellectual state which transcends the world and is free from all material propensity. 
\end{itemize}
When we are in this state, the Logos of God enters into us, putting an end to our battle with the passions and the
demons. [\#94] 

\end{quotex}


\flrightit{Posted on 2014-10-05 by Cologero }

\begin{center}* * *\end{center}

\begin{footnotesize}\begin{sffamily}



\texttt{scardanelli on 2014-10-08 at 14:58 said: }

``More recently, Hans Urs von Balthazar and Valentin Tomberg are representatives of the Church of John."

For whoever might be interested, I came across an interesting article by a fellow named Kevin Mongrain entitled
Rule-Governed Christian Gnosis: Hans Urs von Balthasar on Valentin Tomberg's Meditations on the
Tarot. It's actually quite good for a scholarly article. IT discusses
Tomberg's aims in writing MOT and how these align with the work of von Balthazar.

[Admin]The referenced text is found here: \url{https://www.scribd.com/doc/238193631/Kevin-Mongrain-Rule-governed-Christian-Gnosis}


\hfill

\texttt{Bill McEnaney on 2014-10-13 at 13:23 said: }

Catholics distinguish among kinds s Heaven that ``kingdom of Heaven." signifies. For example, represents the Holy
Trinity's presence in a person's soul. In another, it represents the
condition a disembodied soul is in when it's seeing God face to face. In still another another
one, it's about a physical place that Robert J. Siscoe describes in this article.

\url{http://www.christianorder.com/features/features\_2012/features\_may12.html}


\hfill

\texttt{Fish on 2014-10-28 at 19:34 said: }

Is this Church of John related to the lost consolamentum and the baptism of fire of the medieval christian heretics?


\hfill

\texttt{Cologero on 2014-10-28 at 20:19 said: }

Not at all, Mr Fish, since there is just one Baptism. The reference is to John 21:20-25


\hfill

\texttt{fish on 2014-10-28 at 21:32 said: }

The Gospels speak of a baptism by fire beyond that of water. This double baptism and double regeneration corresponds to
the two phases of initiatic work called albedo and rubedo in Hermetism.


\hfill

\texttt{Cologero on 2014-11-05 at 22:56 said: }

Good point, fish, but that is not a sacrament that someone can do for you.


\end{sffamily}\end{footnotesize}
